\definecolor{hl1}{RGB}{51,51,51}
\definecolor{hl2}{RGB}{24,54,145}
\definecolor{hl3}{RGB}{0,134,179}
\definecolor{hl4}{RGB}{15,0,0}
\definecolor{hl5}{RGB}{3,3,3}
\definecolor{hl6}{RGB}{99,163,92}
\definecolor{hl7}{RGB}{121,93,163}
\definecolor{hl8}{RGB}{167,29,93}
\definecolor{hl9}{RGB}{223,80,0}
\definecolor{hl10}{RGB}{150,152,150}
Code chunks that produce plots automatically generate figures. Use \texttt{fig-cap} for a caption and a label in the chunk header for cross-referencing.




\begin{srccode}
\begin{Verbatim}[commandchars=\\\{\}]
\textcolor{hl1}{```\{r, fig}\textcolor{hl8}{-}\textcolor{hl1}{scatter\}}
\textcolor{hl10}{\#| echo = fenced
}
\textcolor{hl10}{\#| fig-cap = "A simple scatter plot"
}
\textcolor{hl3}{plot}\textcolor{hl1}{(}
\textcolor{hl1}{    }\textcolor{hl3}{1}\textcolor{hl8}{:}\textcolor{hl3}{10}\textcolor{hl1}{,}
\textcolor{hl1}{    (}\textcolor{hl3}{1}\textcolor{hl8}{:}\textcolor{hl3}{10}\textcolor{hl1}{)}\textcolor{hl8}{\^{}}\textcolor{hl3}{2}\textcolor{hl1}{,}
\textcolor{hl1}{    }\textcolor{hl3}{main}\textcolor{hl1}{ }\textcolor{hl8}{=}\textcolor{hl1}{ }\textcolor{hl2}{"Square Numbers"}\textcolor{hl1}{,}
\textcolor{hl1}{    }\textcolor{hl3}{xlab}\textcolor{hl1}{ }\textcolor{hl8}{=}\textcolor{hl1}{ }\textcolor{hl2}{"x"}\textcolor{hl1}{,}
\textcolor{hl1}{    }\textcolor{hl3}{ylab}\textcolor{hl1}{ }\textcolor{hl8}{=}\textcolor{hl1}{ }\textcolor{hl2}{"x\^{}2"}
\textcolor{hl1}{)}
\textcolor{hl1}{```}
\end{Verbatim}
\end{srccode}

\begin{figure}
\centering
\includegraphics[width=0.70\textwidth]{authoring/figures_files/fig-scatter-1.pdf}
\caption{A simple scatter plot}
\label{fig-scatter}
\end{figure}

Figures with labels can be cross-referenced with \texttt{Figure \ref{fig-label}} (see \hyperref[fig-scatter]{Figure 1} above).

Any plotting library works. Here is a ggplot2 example:




\begin{srccode}
\begin{Verbatim}[commandchars=\\\{\}]
\textcolor{hl1}{```\{r, fig}\textcolor{hl8}{-}\textcolor{hl1}{ggplot\}}
\textcolor{hl10}{\#| echo = fenced
}
\textcolor{hl10}{\#| fig-cap = "Fuel efficiency by weight (ggplot2)"
}
\textcolor{hl3}{library}\textcolor{hl1}{(ggplot2)}
\textcolor{hl3}{ggplot}\textcolor{hl1}{(mtcars, }\textcolor{hl3}{aes}\textcolor{hl1}{(wt, mpg, }\textcolor{hl3}{color}\textcolor{hl1}{ }\textcolor{hl8}{=}\textcolor{hl1}{ }\textcolor{hl3}{factor}\textcolor{hl1}{(cyl))) }\textcolor{hl8}{+}
\textcolor{hl1}{  }\textcolor{hl3}{geom\_point}\textcolor{hl1}{(}\textcolor{hl3}{size}\textcolor{hl1}{ }\textcolor{hl8}{=}\textcolor{hl1}{ }\textcolor{hl3}{3}\textcolor{hl1}{) }\textcolor{hl8}{+}
\textcolor{hl1}{  }\textcolor{hl3}{labs}\textcolor{hl1}{(}\textcolor{hl3}{x}\textcolor{hl1}{ }\textcolor{hl8}{=}\textcolor{hl1}{ }\textcolor{hl2}{"Weight (1000 lbs)"}\textcolor{hl1}{, }\textcolor{hl3}{y}\textcolor{hl1}{ }\textcolor{hl8}{=}\textcolor{hl1}{ }\textcolor{hl2}{"MPG"}\textcolor{hl1}{, }\textcolor{hl3}{color}\textcolor{hl1}{ }\textcolor{hl8}{=}\textcolor{hl1}{ }\textcolor{hl2}{"Cylinders"}\textcolor{hl1}{) }\textcolor{hl8}{+}
\textcolor{hl1}{  }\textcolor{hl3}{theme\_minimal}\textcolor{hl1}{()}
\textcolor{hl1}{```}
\end{Verbatim}
\end{srccode}

\begin{figure}
\centering
\includegraphics[width=0.70\textwidth]{authoring/figures_files/fig-ggplot-1.pdf}
\caption{Fuel efficiency by weight (ggplot2)}
\label{fig-ggplot}
\end{figure}

\subsection*{Size}
\label{size}

In Quarto, figure appearance is governed by four parameters: \texttt{fig-width}, \texttt{fig-height}, \texttt{fig-asp}, and \texttt{out-width}. Unless users pick exactly the right combination, figures end up with inconsistent font sizes, axis labels, and line widths. A narrow figure has tiny text while a wide figure has huge text, because the graphics device size changes but the display size does not (or vice versa).

\emph{Calepin} uses pragmatic defaults and adjusts some parameters automatically to maintain a consistent look. In practice, we normally only need two options:

\medskip
\begin{tblr}{colspec={Q[l]Q[l]X[l]}, hline{1,2,Z}={solid}}
Option & Default & What it controls \\
\texttt{out-width} & \texttt{70\%} & Horizontal space in the document \\
\texttt{fig-asp} & \texttt{0.618} & Aspect ratio (golden ratio) \\
\end{tblr}
\medskip

\emph{Calepin} automatically scales the internal graphics device (\texttt{fig-width}) to match \texttt{out-width}, so text and labels inside figures stay the same size regardless of how wide the figure is displayed. This means a narrow figure and a wide figure will have consistent font sizes.

\begin{itemize}
\item Want a wider figure? Change \texttt{out-width}.


\item Want a taller figure? Change \texttt{fig-asp}.


\end{itemize}

That is all.




\begin{srccode}
\begin{Verbatim}[commandchars=\\\{\}]
\textcolor{hl1}{```\{r, chunk}\textcolor{hl8}{-}\textcolor{hl3}{3}\textcolor{hl1}{\}}
\textcolor{hl10}{\#| fig-asp = 0.4
}
\textcolor{hl10}{\#| out-width = "95\%"
}
\textcolor{hl10}{\#| fig-cap = "A wide, short figure"
}
\textcolor{hl10}{\#| echo = fenced
}
\textcolor{hl3}{plot}\textcolor{hl1}{(}\textcolor{hl3}{1}\textcolor{hl8}{:}\textcolor{hl3}{20}\textcolor{hl1}{, }\textcolor{hl3}{sin}\textcolor{hl1}{(}\textcolor{hl3}{1}\textcolor{hl8}{:}\textcolor{hl3}{20}\textcolor{hl1}{), }\textcolor{hl3}{type}\textcolor{hl1}{ }\textcolor{hl8}{=}\textcolor{hl1}{ }\textcolor{hl2}{"l"}\textcolor{hl1}{, }\textcolor{hl3}{col}\textcolor{hl1}{ }\textcolor{hl8}{=}\textcolor{hl1}{ }\textcolor{hl2}{"steelblue"}\textcolor{hl1}{, }\textcolor{hl3}{lwd}\textcolor{hl1}{ }\textcolor{hl8}{=}\textcolor{hl1}{ }\textcolor{hl3}{2}\textcolor{hl1}{)}
\textcolor{hl1}{```}
\end{Verbatim}
\end{srccode}

\begin{figure}
\centering
\includegraphics[width=0.95\textwidth]{authoring/figures_files/chunk-3-1.pdf}
\caption{A wide, short figure}
\label{chunk-3}
\end{figure}

The \texttt{fig-width} and \texttt{fig-height} options still exist for full manual control, but you rarely need them. When you set \texttt{fig-width} explicitly, auto-scaling is disabled.

\subsection*{Options}
\label{options}

\medskip
\begin{tblr}{colspec={Q[l]Q[l]X[l]}, hline{1,2,Z}={solid}}
Option & Default & Description \\
\texttt{out-width} & \texttt{70\%} & Display width in the document \\
\texttt{fig-asp} & \texttt{0.618} & Aspect ratio (height / width) \\
\texttt{fig-width} & \texttt{6} & Graphics device width in inches (auto-scaled from \texttt{out-width}) \\
\texttt{fig-height} &  & Derived from \texttt{fig-width * fig-asp} \\
\texttt{fig-cap} &  & Caption text \\
\texttt{fig-alt} &  & Alt text for accessibility \\
\texttt{fig-align} & \texttt{center} & Alignment (\texttt{left}, \texttt{center}, \texttt{right}) \\
\texttt{dev} & \texttt{png} & Graphics device (\texttt{png}, \texttt{pdf}, \texttt{svg}) \\
\end{tblr}
\medskip

\subsection*{Layouts}
\label{layouts}

Figures can be arranged in grids using \texttt{layout-ncol}. The following example places three plots side by side with individual subcaptions and an overall caption.



\begin{figure}
\centering

\begin{minipage}[t]{0.320\textwidth}
\centering
\includegraphics[width=1.00\textwidth]{authoring/figures_files/line-1.pdf}
\caption{Linear}
\label{line}
\end{minipage}
\hfill
\begin{minipage}[t]{0.320\textwidth}
\centering
\includegraphics[width=1.00\textwidth]{authoring/figures_files/quad-1.pdf}
\caption{Quadratic}
\label{quad}
\end{minipage}
\hfill
\begin{minipage}[t]{0.320\textwidth}
\centering
\includegraphics[width=1.00\textwidth]{authoring/figures_files/log-1.pdf}
\caption{Logarithmic}
\label{log}
\end{minipage}
\caption{Three growth patterns compared.}
\label{fig-panels}
\end{figure}

See \hyperref[fig-panels]{Figure 3} for the result. The \texttt{layout} attribute supports custom grids with relative widths (e.g., \texttt{layout=\textquotedbl{}{[}{[}1,1{]},{[}1{]}{]}\textquotedbl{}}) and negative values for spacing gaps.

